% q0018.tex — Going with the Flow (Steady-State Unemployment)
\question{
  id        = {0018},
  title     = {Going with the Flow},
  source    = {OBECON},
  year      = {2026},
  round     = {Seletiva IEO -- Phase A},
  language  = {English},
  topic     = {Macro},
  subtopic  = {Labor Markets / Steady-State Unemployment},
  type      = {objective},
  statement = {%
    In a labor market, let $f$ be the \emph{job-finding rate} (probability an
    unemployed person finds a job) and $s$ be the \emph{job separation rate}
    (probability an employed person loses their job). Both rates are constant.

    What is the steady-state unemployment rate?
  },
  choices   = {%
    \item $u^* = \dfrac{f}{f+s}$
    \item $u^* = \dfrac{s}{s+f}$
    \item $u^* = \dfrac{s-f}{s+f}$
    \item $u^* = s \cdot f$
    \item $u^* = \dfrac{1}{s+f}$
  },
  answer    = {B},
  solution  = {%
    At steady state, flows \emph{into} unemployment equal flows \emph{out}
    of unemployment:
    \[
      \underbrace{s \cdot (1 - u^*)}_{\text{employed} \to \text{unemployed}}
      = \underbrace{f \cdot u^*}_{\text{unemployed} \to \text{employed}}
    \]
    Solving for $u^*$:
    \[
      s - s u^* = f u^* \implies s = u^*(f + s)
      \implies \boxed{u^* = \frac{s}{s + f}}
    \]

    Intuition: a higher separation rate $s$ raises unemployment; a higher
    job-finding rate $f$ lowers it. Answer: \textbf{(B)}.
  },
}
